\documentclass[11pt]{article}
\usepackage[letterpaper,margin=0.88in]{geometry}
\usepackage[T1]{fontenc}
\usepackage{lmodern,microtype}
\usepackage{amsmath,amssymb,amsthm,mathtools}
\usepackage{booktabs,longtable,array}
\usepackage{xcolor,enumitem,fancyhdr,needspace}
\usepackage[hyphens]{url}
\usepackage[colorlinks=true,linkcolor=blue!45!black,urlcolor=blue!55!black,citecolor=blue!45!black]{hyperref}
\usepackage{listings}
\lstset{basicstyle=\ttfamily\small,breaklines=true,columns=fullflexible,frame=single,
 rulecolor=\color{black!15},backgroundcolor=\color{black!2},showstringspaces=false}
\setlength{\parindent}{0pt}
\setlength{\parskip}{6pt}
\setlength{\emergencystretch}{2em}
\setlist{itemsep=3pt,topsep=4pt}
\newtheorem{theorem}{Theorem}[section]
\newtheorem{lemma}[theorem]{Lemma}
\newtheorem{proposition}[theorem]{Proposition}
\newtheorem{corollary}[theorem]{Corollary}
\theoremstyle{remark}
\newtheorem{remark}[theorem]{Remark}
\newcommand{\Z}{\mathbb Z}
\newcommand{\N}{\mathbb N}
\newcommand{\Pos}{\mathbb Z_{>0}}
\newcommand{\Con}{\operatorname{Con}(\mathrm{ZFC})}
\newcommand{\sha}{\texttt{84e16e43092614bae262ca59651cb9d77e31a255}}
\newcommand{\fl}[1]{\left\lfloor#1\right\rfloor}
\newcommand{\rg}[2]{\{#1,\ldots,#2\}}
\pagestyle{fancy}
\fancyhf{}
\fancyhead[L]{\small Diophantine consistency equations}
\fancyhead[R]{\small 24 October 2026}
\fancyfoot[C]{\thepage}
\setlength{\headheight}{14pt}
\hypersetup{pdftitle={Removing Redundancy from Diophantine Consistency Equations},
 pdfauthor={AI-assisted research prepared for Vladimir Reshetnikov},
 pdfsubject={Explicit equations, proofs, and exact computational verification}}
\begin{document}
\hypersetup{pageanchor=false}
\begin{titlepage}
\vspace*{0.3in}
{\small\scshape Research for the ProveIt project\par}
\vspace{0.35in}
{\huge\bfseries\raggedright Removing Redundancy from\\
Diophantine Consistency Equations\par}
\vspace{0.25in}
{\Large Unweighted Pell radicals, stronger transport bounds,\\
shared integer guards, and exact norm identities\par}
\vspace{0.4in}
{\large Prepared for Vladimir Reshetnikov\\24 October 2026\par}
\vspace{0.3in}
\begin{abstract}
We construct and verify several improvements to explicit polynomial equations
whose unsolvability is equivalent to syntactic consistency of ZFC.
A direct quartic over the unrestricted integers is reduced from 604 to
589 unknowns and from 9,676 to 8,156 monomials, at unchanged coefficient
height. A square-gap theorem removes an asymmetric radical weight, and an
auxiliary Pell index bound eliminates a separate transport majorant.
Together they give a ten-positive-variable equation with 96 arithmetic
operations and exact degree 360, and a direct-base companion with
102 operations and exact degree 332.
Exact norm identities give pointwise-identical integer evaluators with
126 operations in thirteen variables, and 138 or 137 in twelve variables.

Complete equations, fixed coefficient payloads, source pins, builders and
exact checkers accompany the report. All comparisons use one frozen
repository snapshot. These results improve specific complexity measures;
they do not establish global minimality or reduce the repository's best
variable count. The arithmetic-to-refutation calibration is inherited.
No new end-to-end Lean, Rocq or object-language ZFC derivation is claimed.
\end{abstract}
\vfill
{\small AI-assisted research. Unbounded mathematical arguments are stated
separately from exact algebraic and finite-specification checks.}
\end{titlepage}
\pagenumbering{roman}
\hypersetup{pageanchor=true}
\begingroup
\small\setlength{\parskip}{0pt}
\tableofcontents
\endgroup
\clearpage
\pagenumbering{arabic}

\section{Results, domains and the meaning of simplicity}

The objective is a fixed polynomial $P$ with integer coefficients and a
precisely stated domain $D$ for which ordinary mathematics establishes
\[
 \mathrm{ZFC}\vdash
 \left[(\forall\boldsymbol{x}\in D^n)\,
        P(\boldsymbol{x})\ne0\right]\ \longleftrightarrow\ \Con.
 \tag{1}
\]
Here $\Con$ means that the specified syntactic ZFC calculus has no finite
derivation of falsity. We neither assume nor prove $\Con$.
A proof of the biconditional is different from a proof of its left side.

There is no established scalar measure of the \emph{simplest} Diophantine
equation. Variable count, degree, arithmetic operations, coefficient height
and expanded support measure different costs. We preserve the domain and
compare those quantities separately. Ten positive variables are not ten
unrestricted integer variables.

\begin{table}[ht]
\centering\small
\begin{tabular}{@{}lrrrr@{}}
\toprule
Domain and construction & Variables & Old degree & New degree & Operations\\
\midrule
Positive, scaled base &10&368&\textbf{360}&99 $\to$ \textbf{96}\\
Positive, direct base &10&340&\textbf{332}&105 $\to$ \textbf{102}\\
Integers, quadratic trace &13&536&536&135 $\to$ \textbf{126}\\
Integers, rational carrier &12&1,200&1,200&145 $\to$ \textbf{138}\\
Integers, integral carrier &12&1,456&1,456&141 $\to$ \textbf{137}\\
\bottomrule
\end{tabular}
\caption{Exact circuit comparisons. Each live binary addition, subtraction
or multiplication costs one operation, including multiplication by a
constant. Constant leaves and register reuse are free.}
\label{tab:circuits}
\end{table}

The direct quartic gives a different improvement:
\[
\begin{array}{c|rr}
 &\text{parent}&\text{new}\\ \hline
\text{unrestricted integer variables}&604&589\\
\text{quadratic residuals}&372&371\\
\text{total degree}&4&4\\
\text{collected monomials}&9,676&8,156\\
\text{maximum absolute coefficient}&
2,445,790,974,704&2,445,790,974,704
\end{array}
\tag{2}
\]
The saving is fifteen actual variables after projection, and 1,520
monomials, approximately $15.71\%$ of the parent's support.
All 8,156 coefficients are supplied explicitly.

The comparison repository is \texttt{VladimirReshetnikov/ProveIt}.
The frozen commit is
\begin{center}\small\sha.\end{center}
Its active research is newer than its placed manuscript
collection. The 99/368 and 105/340 positive endpoints, thirteen-integer
trace, two twelve-integer carriers, and 75-predicate quartic family are
the actual parents used here \cite{zero,direct,trace,scalar,quartic75}.

\subsection{Different kinds of equivalence}

The integer norm rewrites are \emph{identities of polynomials}: values and
zero sets agree at every assignment, including degenerate ones.
Removing the radical weight and simplifying the transport carrier changes
the positive polynomial; we prove equivalence of \emph{positive zero
existence}, not equality of zero sets at unchanged inputs.
The quartic transformations change certificate variables and reconstruct
every old sign condition before invoking the source interpretation.

The contributions are new improvements relative to the inspected
snapshot. We do not claim worldwide priority for a generic norm identity
or an elementary sign lemma. The finite repository audit is not an
exhaustive historical search.

\subsection{Inherited consistency calibration}

The positive equations and the small integer norm equations inherit the
repository's ordinary native-compiler/K215/refutation calibration.
The direct quartic uses its separate C222 proof-certificate route and
does not depend on K215 or the Jones nine-variable compression.
The two routes are not identified.

The new finite algebraic transformations and ordinary witness-existence
arguments are formalizable in ZFC. Composing them with their pinned
source theorems gives (1). Exact checkers do not replace the inherited
unbounded calibration proof, and this report does not supply the complete
encoded derivation in the project's object-language ZFC calculus.

Small-variable universal constructions are separate substantial results;
Sun's ten-integer theorem illustrates the distinction between reducing
variable count and simplifying a particular evaluator \cite{sun}.
The present packet does not improve that general bound or the repository's
nine-positive-variable endpoint.

\section{Exact four-root norm identities and smaller integer evaluators}
\label{sec:norm}

For formal radicands $a,b,c,d$, define
\[
F_{16}(x;a,b,c,d)=
\prod_{\epsilon\in\{-1,1\}^4}
\bigl(x-\epsilon_1\sqrt a-\epsilon_2\sqrt b
        -\epsilon_3\sqrt c-\epsilon_4\sqrt d\bigr).
\tag{3}
\]
Changing the sign of a root permutes the factors. The product is therefore
an integer polynomial, even in $x$. Write
$N_4(z;a,b,c,d)=F_{16}(x;a,b,c,d)$ when $z=x^2$.

\subsection{A symmetric identity}

Compute
\[
\begin{aligned}
p&=a+b,&q&=c+d,&u&=ab,&v&=cd,\\
e_1&=p+q,&e_2&=u+v+pq,&e_3&=uq+vp,&e_4&=uv,
\end{aligned}
\]
\[
s=z-e_1,\qquad t=s^2-4e_2,\qquad
C=t^2+64(e_4-ze_3),\qquad L=t-4zs.
\tag{4}
\]

\begin{theorem}
The following is an identity of integer polynomials:
\[
\boxed{N_4(z;a,b,c,d)=C^2-256e_4L^2.}
\tag{5}
\]
The displayed schedule uses 29 binary operations when $z$ is supplied.
\end{theorem}

\begin{proof}
For algebraically independent $r_1,r_2,r_3,r_4$, put $a_i=r_i^2$ and
$x=r_1+r_2+r_3+r_4$. Let $p_2,p_3,p_4$ be the elementary symmetric
expressions in the $r_i$. Direct expansion gives
\[
2p_2=x^2-e_1=s,\quad e_2=p_2^2-2xp_3+2p_4,\quad
e_3=p_3^2-2p_2p_4,\quad e_4=p_4^2.
\]
Consequently $t+8p_4=8xp_3$ and $p_3^2=e_3+sp_4$. Squaring and
substituting yields
\[
t^2+64(e_4-x^2e_3)+16p_4(t-4x^2s)=0.
\]
Thus $C+16p_4L=0$; eliminating $p_4$ proves that (5) vanishes at
$x=\sum r_i$. The calculation applies to all sixteen sign choices.
Those roots are distinct over the generic rational-function field.
Both expressions are monic of degree sixteen in $x$, so they are equal.
Polynomial specialization covers zero radicands and coincident roots;
no division by a specialized quantity occurs.

The four elementary symmetric expressions cost twelve operations.
The remaining schedule costs $1+3+5+3+5=17$, giving 29.
\end{proof}

\subsection{The preferred 28-operation schedule}

During review, a retained source proof was found to contain a 28-operation
schedule for the same norm. It is explicitly credited and pinned in the
packet. It gives the preferred final counts, one operation below the
symmetric variants:
\[
\begin{aligned}
\sigma_1&=a+b+c,&\sigma_2&=ab+(a+b)c,&\sigma_3&=abc,\\
U&=z+d,&P&=U-\sigma_1,&
E&=P^2+4(zd-\sigma_2),\\
G_{\rm e}&=E^2+16(zdP^2-4\sigma_3U),&
G_{\rm o}&=EP-16\sigma_3,\\
N_4(z;a,b,c,d)&=G_{\rm e}^2-64zdG_{\rm o}^2.
\end{aligned}
\tag{6}
\]
The cost is 28 with shared intermediates. Pair the two choices of the
fourth root in the octic three-root norm to derive (6), or compare all
coefficients with (3). The delivered checker uses the latter independent
method, obtaining all 495 coefficients.

The main computational improvement is to clear a rational radicand
\emph{before} applying either compact norm schedule:
\[
\boxed{
\delta^8F_{16}(Z;A,B,E,C_{\rm n}/\delta)
=N_4(\delta Z^2;\delta A,\delta B,\delta E,C_{\rm n}).
}
\tag{7}
\]
The left side means its cleared polynomial. Homogeneity proves equality
for nonzero $\delta$, and polynomial identity extends it to $\delta=0$.
No actual division is required in the evaluator.

If $A=U^2D_f$, $B=U^2D_t$, $E=U^2A_g$, and $C_{\rm n}=U^2C_y$,
form $D=\delta U^2$ once and evaluate
\[
N_4(\delta Z^2;DD_f,DD_t,DA_g,U^2C_y).
\tag{8}
\]
The shared scale saves two additional multiplications. With (6), the
whole thirteen-integer evaluator drops from 135 to 126 operations.
The alternative symmetric evaluator, with the elementary identity (5),
has 127 operations and is also retained.

\subsection{A complete thirteen-integer equation}

The variables are
\[
j,C,z,h,r,v,w,T,c,n,t,p_2,p_3\in\Z.
\tag{9}
\]
The six fixed coefficients
$\beta,K_0,M_C,M_F,A_{\rm cone},B_{\rm cone}$ are the entire pinned
mixed-parity tuple in \texttt{integer/coefficients/}. In particular,
\[
\beta=2^{1064921}-1,\qquad A_{\rm cone}=2^{1113},\qquad
B_{\rm cone}=2^{666104}.
\]
The other three numerals are supplied as complete, authenticated
little-endian byte strings. These are fixed coefficients, not quantified
inputs. This tuple differs from the BFS tuple of Appendix~\ref{app:constants}.
The supplied $M_F$ already includes its single shift by $\beta$.

Define
\[
\begin{aligned}
J&=6j+1,\quad q=\beta J+1,\quad W=12w+4,\quad X=qW,\\
Y&=qv,\quad E_0=XY,\quad a=E_0+Y,\quad H=4a+3,\\
\Delta&=a^2+H,\qquad F=(K_0+W)C-z(q-1),\\
R&=[q(q-F)-C](q^2-1)+(M_C+qM_F)J,\\
k&=R+1+hE_0,\quad\eta=c-Yk,\quad\zeta=k-\eta,\\
i&=2r+1,\quad I=ic^2,\quad D_f=1+\Delta I^2,\quad Q=\Delta^2I^2,\\
\delta&=Q-1,\quad V=D_f(cT-R)-c,\quad C_y=QV^2-1,\\
\ell&=XY^2,\quad L_0=\ell k,\quad D_t=1+L_0(L_0+k).
\end{aligned}
\tag{10}
\]
The remaining expressions are
\[
\begin{aligned}
G&=(A_{\rm cone}F)(B_{\rm cone}C)
 -(q-A_{\rm cone}F-B_{\rm cone}C)^2,\\
\Pi&=G\eta\zeta-(t^2-G-\eta\zeta+1)^2,\\
A_g&=\Pi-1-p_2^2-p_3(p_3+1),\\
M&=2q^2,\quad U=MH,\quad K=X(X+2ac)-1,\\
Z&=U(n+V)+MK+H(v-q^2),\qquad D=\delta U^2.
\end{aligned}
\tag{11}
\]
The equation is (8) equal to zero, with the explicit norm (6).
It is an ordinary polynomial in precisely the thirteen variables (9).
Identity (7) proves equality with the entire parent polynomial.
Its exact degree is 536; for example the attained-degree coefficient is
\[
[j^{240}w^{80}v^{96}h^{16}r^{32}c^{64}z^8]P
=2^{488}3^{320}\beta^{240}>0.
\tag{12}
\]
The exponents sum to 536. Universal polynomial equality is why this
coefficient, the degree, and every integer zero transfer.

\subsection{Twelve-integer companions}

The two companions use $j,C,z,h,r,v,w,T,c,n,t,s\in\Z$ and the same
six coefficients. Their complete prefixes are unchanged from the pinned
scalar-carrier sources. In the source notation, their new tails are
\[
N_4(\delta Z^2;
9\delta S^2D_f,\ 81\delta S^2D_t,\ 729\delta S^2D_g,\ S^2C_y)
\tag{13}
\]
and
\[
N_4(Z^2;9S^2D_f,\ 81S^2D_t,\ 729S^2D_g,\ S^2\delta C_y).
\tag{14}
\]
With (6) they have 138 operations at degree 1,200 and 137 at degree 1,456.
The symmetric schedule (5) costs one additional operation in each case.
The complete definitions of $S,Z,D_g$ and every other front abbreviation
are literal in the delivered JSON circuits and copied source proofs.
These are exactly the original polynomials, with smaller evaluators.

\section{The new ten-positive-variable equation}
\label{sec:positive}

The only unknowns are
\[
J,F,\alpha,i,T,w,Y,n,\eta,\zeta\in\Pos.
\tag{15}
\]
Use the four BFS coefficients $\beta,K_0,M_C,M_F$ defined entirely in
Appendix~\ref{app:constants}. All other symbols below are polynomial
abbreviations.
\[
\begin{aligned}
q&=\beta J+1,&C&=q-F-\alpha,&X&=wq,\\
k&=\eta+\zeta,&c&=kY+\eta,&\mu&=q^2+1,\\
t&=\mu Y,&A&=t(X+1)+1,&\Delta&=A^2-1,\\
H&=4A-5,&v_0&=Xt^2,&I&=ic^2.
\end{aligned}
\tag{16}
\]
Set
\[
\begin{aligned}
D_\tau&=1+v_0(v_0+1)k^2,&D_f&=1+\Delta I^2,\\
Q&=\Delta^2I^2,&\delta&=Q-1,\\
R&=[(q-1)(q-F)+\alpha](q^2-1)+(M_C+qM_F)J,\\
e&=k-R-1,&B_t&=(K_0+w)C-F,\\
b&=c(T+\alpha)-R,&V&=D_fb-c,\\
D_{\rm n}&=QV^2-1,&K_{\rm pos}&=X(3c-2X)+2,\\
W&=1+(q-1)b,&P_0&=q^3(q-1),\qquad M=P_0\mu^2H.
\end{aligned}
\tag{17}
\]
Form
\[
\begin{split}
Z={}&P_0\mu^2[H(n+V)+K_{\rm pos}]\\
&+H[\mu^2((q-1)Y-q^3WB_t)-P_0e],
\end{split}
\tag{18}
\]
\[
S=\delta Z^2,\quad B_\tau=\delta M^2D_\tau,\quad
B_f=\delta M^2D_f,\quad B_y=M^2D_{\rm n},
\tag{19}
\]
\[
E=(S+B_y-B_\tau-B_f)^2-4(SB_y+B_\tau B_f).
\tag{20}
\]
The single equation is
\[
\boxed{\qquad P_{96}=E^2-64SB_yB_\tau B_f=0.\qquad}
\tag{21}
\]
It has exact degree 360 and a 96-operation evaluator.
There is no square root, division, variable exponent, hidden coefficient,
or additional quantified variable in (16)--(21).
The circuit evaluates $D_\tau$ as $(v_0k)(v_0k+k)+1$ and shares
$A-1=(\mu Y)(X+1)$.

Relative to the pinned 99/368 parent, two changes are made:
the factor 9 on $B_\tau$ is removed, and the transport factor
$1+(q-1)(HI+b)$ is replaced by $1+(q-1)b$.
The first requires separation of three radicals; the second requires a
stronger transport bound. Neither change is justified by an identity of
the complete old and new polynomials.

The arithmetic argument is uniform for coherent tuples satisfying
\[
\beta\text{ odd},\quad\beta\ge31,\quad K_0>0,\quad
0<M_C<\beta<M_F<2\beta.
\tag{22}
\]
These inequalities are not themselves a consistency calibration.
All four actual constants must be the complete calibrated tuple, as they
are in the appendix and the reconstructed coefficient files.

\section{A square gap that removes the radical weight}
\label{sec:gap}

The parent uses $9B_\tau$ in place of $B_\tau$, so its signed radical
relation contains $3\tau$. Deleting that weight requires more than checking
the intended integral witnesses: a rational signed sum of irrational
roots could otherwise introduce extra solutions. The following argument
excludes every such cancellation before any source semantics is used.

\begin{lemma}[A square gap]
\label{lem:gap}
Let $x\ge2$ and $t\ge1$ be integers. Put
\[
\Delta=t(x+1)(t(x+1)+2),\qquad v=t^2x,
\quad D=xt(x+1)(t(x+1)+2)(t^2x+1).
\]
Then $\Delta v(v+1)=t^2D$ is not a rational square.
In particular $\Delta I^2=v(v+1)k^2$ has no solution with nonzero
integers $I,k$.
\end{lemma}

\begin{proof}
Set $m=\lfloor x/2\rfloor$ and $h=t^2x(x+1)+tx+m$.
For $x=2m$, direct expansion gives
\[
D-h^2=m[2t^2+4(m+1)t-m]>0,
\tag{23}
\]
\[
(h+1)^2-D=2m(4m+1)t^2-4m^2t+(m+1)^2>0.
\tag{24}
\]
Here $m\ge1$. The first is positive already after replacing $t$ by 1.
In the second, $t^2\ge t$ gives the positive lower bound
$(4m^2+2m)t+(m+1)^2$.
For $x=2m+1$, where $m\ge1$, the exact identities are
\[
D-h^2=t^2x(x+2)+2tx(m+2)-m^2>0,
\tag{25}
\]
\[
(h+1)^2-D=[tx-(m+1)]^2>0.
\tag{26}
\]
The first is positive at $t=1$ and increases thereafter. In the second,
$tx\ge2m+1>m+1$. Thus $h^2<D<(h+1)^2$ in both parity classes.
An integer that is a rational square is an integer square, by reduction
of its rational square root to lowest terms.
Finally the prohibited equation would imply $D=(\Delta I/(tk))^2$.
The excluded point $x=t=1$ gives $D=16$; the domain is essential.
\end{proof}

\begin{lemma}[Rationality of three separated radicals]
\label{lem:rational}
Let $a,b,c$ be positive, pairwise distinct rational numbers with
\[
(a+b-c)^2\ne4ab.
\tag{27}
\]
If $u^2=a$, $v^2=b$, $w^2=c$ and $u+v+w$ is rational, then
$u,v,w$ are rational.
\end{lemma}

\begin{proof}
Write $s=u+v+w$. If $s=0$, squaring $w=-u-v$ twice contradicts
(27), so $s\ne0$. For $K_a=s^2+a-b-c$, squaring
$2vw=K_a-2su$ gives
\[
4sK_a u=K_a^2+4s^2a-4bc.
\tag{28}
\]
At least one of $K_a,K_b,K_c$ is nonzero, since otherwise subtracting
them would give $a=b=c$. The corresponding identity makes one root
rational. The sum of the other two is rational and nonzero, since their
squares differ. If that sum is $r=v+w$, then
$v=(r^2+b-c)/(2r)$ and $w=r-v$.
\end{proof}

\begin{theorem}[Unweighted separation]
\label{thm:separation}
Use $\Delta,v$ from Lemma~\ref{lem:gap}. Suppose an integer $\mu>3$
divides $t$, integers $I,k$ satisfy $I\ge k\ge1$, and $V$ is a
nonzero integer. Put
\[
d_\tau=1+v(v+1)k^2,\quad d_f=1+\Delta I^2,\quad
Q=\Delta^2I^2,\quad d_y=\frac{QV^2-1}{Q-1}.
\tag{29}
\]
Every rational signed sum of $\sqrt{d_\tau},\sqrt{d_f},\sqrt{d_y}$
has all three radicals rational.
\end{theorem}

\begin{proof}
We have $\Delta>v+1>1$, so $Q>1$ and all three radicands are positive.
Lemma~\ref{lem:gap} excludes $d_\tau=d_f$.
If $d_\tau=d_y$, then $Q(d_\tau-V^2)=d_\tau-1$,
but its right side is strictly between 0 and $Q$, since
$v(v+1)k^2<\Delta^2I^2$.
If $d_f=d_y$, rearrangement gives $d_f-V^2=1/\Delta$,
contradicting integrality. Thus they are pairwise distinct.

For the remaining triangle obstruction, put $\delta=Q-1$ and
$D_{\rm n}=QV^2-1$. Equality in (27) would imply
\[
[\delta(d_\tau+d_f)-D_{\rm n}]^2-4\delta^2d_\tau d_f=0.
\tag{30}
\]
Modulo $\mu$ we have $\Delta=v=Q=0$,
$\delta=D_{\rm n}=-1$, and $d_\tau=d_f=1$.
The left side is $-3$ modulo $\mu$, impossible for $\mu>3$.
Lemma~\ref{lem:rational} applies, with arbitrary signs on the roots.
\end{proof}

\subsection{Application on the full positive input domain}

In (16), $X\ge32$, $t=\mu Y$, $\mu>3$, $c=kY+\eta>k$ and
$I=ic^2\ge c^2>k$. Also $D_f>2c$ and $b$ is integral.
If $b\le0$, then $V=D_fb-c\le-c$; if $b\ge1$, then
$V\ge D_f-c>c$. Thus $V\ne0$ before the polynomial equation is used.
Every hypothesis of Theorem~\ref{thm:separation} is unconditional on the
supplied positive domain.

The octic identity
\[
F_8(s;a,b,c)=[(s+c-a-b)^2-4(sc+ab)]^2-64scab
\tag{31}
\]
factors, with $s=x^2$, into the eight signed linear radical factors.
Since $M,\delta>0$, equation (21) implies
\[
\frac ZM=\epsilon_1\tau+\epsilon_2 f+\epsilon_3 y,\qquad
\tau^2=D_\tau,\quad f^2=D_f,\quad y^2=D_{\rm n}/\delta,
\tag{32}
\]
with positive roots.
Theorem~\ref{thm:separation} makes all three roots rational.
The first two are integers because their squares are integers; the
retained adjacent-Pell descent lemma makes the third integral under the
unchanged $Q=\Delta^2I^2$ and integer-$V$ hypotheses \cite{zero}.

\section{The auxiliary Pell index supplies its own transport bound}
\label{sec:transport}

Removing the weight alone gives 98 operations at degree 368.
We now show that the already encoded auxiliary Pell solution grows enough
to remove the separate $HI$ term in the transport majorant.
This is the additional reduction to (21).

For $A\ge2$, write
\[
\chi_A(m)+\psi_A(m)\sqrt{A^2-1}
=(A+\sqrt{A^2-1})^m.
\]
The unchanged strong/auxiliary folding argument, after the integral roots
are recovered, gives $c=\psi_A(p)$ for an odd positive $p$.
The parity does not require the later equality $p=R$:
the auxiliary index is odd, and, with $m$ the strong Pell index, both
reduction modulo $2m$ and reflection
from an index $u$ to $2m-u$ preserve oddness.
Since $c\ge3$, $p=1$ is excluded, hence
\[
c\ge\psi_A(3)=4A^2-1=4\Delta+3.
\tag{33}
\]

Before transport, the retained first-index growth and mask estimate give
\[
c>8X\mu^2Y^2,\qquad R<c\alpha/4.
\tag{34}
\]
These bounds do not use $C>0$, the main index equality, or the sign of
$B_t$. In particular
\[
b=c(T+\alpha)-R>c(T+3\alpha/4)>c.
\tag{35}
\]

\begin{lemma}[Auxiliary index at least five]
\label{lem:auxfive}
Under the recovered kernel equations and (33)--(35),
\[
b>8Q>2HI.
\tag{36}
\]
\end{lemma}

\begin{proof}
Set $S_0=\Delta I$, so $Q=S_0^2$.
The auxiliary equation is
$S_0^2V^2-(S_0^2-1)y^2=1$.
Here $V>0$. Pell classification gives
$S_0V=\chi_{S_0}(\nu)$ at a positive index $\nu$.
Even indices have $\chi_{S_0}(\nu)\equiv\pm1\pmod{S_0}$,
whereas $S_0$ divides $S_0V$, so $\nu$ is odd.

By (33)--(35),
\[
V=D_fb-c>\Delta I^2c>4\Delta^2I^2=4Q.
\]
The possibilities $\nu=1$ and $\nu=3$ would give respectively
$V=1$ and $V=4Q-3$, and are impossible.
Therefore $\nu\ge5$, and monotonicity gives
\[
V\ge16Q^2-20Q+5>8Q^2.
\tag{37}
\]
The last inequality holds for $Q\ge3$, far below the actual lower bound.
Because $D_f=\Delta I^2+1<Q$,
$b=(V+c)/D_f>8Q$.
Finally $\Delta-H=(A-2)^2>0$ gives $8Q>2HI$.
\end{proof}

\begin{lemma}[Quarter bounds for the two error terms]
\label{lem:quarter}
With $u=y-V$, the recovered positive roots satisfy
\[
0<u<b/4,\qquad e/\mu^2<b/4.
\tag{38}
\]
\end{lemma}

\begin{proof}
The auxiliary equation gives
\[
u=\frac{V^2-1}{(Q-1)(y+V)}
<\frac{V}{2(Q-1)}<\frac{D_fb}{2(Q-1)}.
\]
Here $y>V>1$, and
$Q-1-2D_f=\Delta(\Delta-2)I^2-3>0$.
Thus $u<b/4$.
Also $b-e=c(T+\alpha)-k+1>0$, so $e<b$.
Since $b>0$ and $\mu^2>4$, $e/\mu^2<b/4$.
This upper bound also holds when $e$ is negative; no premature
divisibility or positivity of $e$ is assumed.
\end{proof}

\subsection{Soundness of the shorter carrier}

The retained bound $3\tau+f<HI$ implies $\tau+f<HI<b/2$.
The normalized center from (18) is
\[
\frac ZM=n+V+\frac{K_{\rm pos}}H+\frac Y{q^3}
-\frac{WB_t}{q-1}-\frac e{\mu^2}.
\tag{39}
\]
Suppose $B_t\le-1$.
With the new $W=1+(q-1)b$, $WB_t/(q-1)<-b$.
The parent growth estimate already gives $K_{\rm pos}>0$.
Using (32), (36) and (38),
\[
\begin{split}
n
&=\epsilon_1\tau+\epsilon_2f+(\epsilon_3y-V)
-K_{\rm pos}/H-Y/q^3+WB_t/(q-1)+e/\mu^2\\
&<b/2+b/4+b/4-b=0,
\end{split}
\]
contrary to $n>0$. Thus $B_t\ge0$ and hence $C>0$.

Everything after this transport step is unchanged.
In the original order the proof recovers $3q<R<q^4-2$, $p=R$, parity,
and then $\mu^2\mid e$. The first index is $(R+1)/2$.
The ratio estimate and whole-product projection force $X=2^R$;
no noninvertible factor is cancelled modulo $H$.
The remaining transport and half-binomial conditions reconstruct the
same unscaled native source and a finite syntactic refutation.
The new $W$ retains residue $1$ modulo $q-1$.

\subsection{Completeness}

Use the parent's fresh canonical zero-offset witnesses from a finite
refutation. Retain every coordinate except $n$. At these witnesses
\[
s_0=Y/q^3\ge1,\quad t_0=B_t/(q-1)\ge0,\quad
\rho=(k-R-1)/\mu^2\ge1,\quad
u=y-V>0,\quad\kappa=K_{\rm pos}/H>0
\]
are integers. With the new $W$, choose
\[
n=Wt_0+\rho+\tau+f+u-s_0-\kappa.
\tag{40}
\]
The parent gives $\tau>k>Y\ge s_0$ and $f>\kappa$.
Also $W>0$ because $b>c>0$, so $n>0$.
Equation (39) becomes $Z/M=\tau+f+y$, and a norm factor vanishes.
This is canonical reconstruction, not an asserted positive pointwise
inverse from every parent tuple.

\begin{corollary}
At the coherent coefficients of Appendix~\ref{app:constants},
$P_{96}$ has a positive integer zero exactly when the pinned ZFC calculus
has a finite refutation. Therefore ordinary mathematics formalizable in
ZFC gives (1) with $D=\Pos$ and $n=10$.
\end{corollary}

\section{Exact degrees and the direct-base companion}
\label{sec:degree}

The degree must be calculated for the entire norm. Following only the
centered difference can miss a different leading term. Write
\[
s=\delta Z^2,\quad a=\delta M^2D_\tau,\quad
b_*=\delta M^2D_f,\quad c_*=M^2D_{\rm n},\quad d=s-c_*.
\]
Then
\[
E=d^2-2(s+c_*)(a+b_*)+(a-b_*)^2.
\tag{41}
\]
For the new scaled carrier the relevant degrees are
\[
\deg a=74,\quad\deg b_*=76,\quad
\deg s=\deg c_*=104,\quad\deg d=86.
\]
The term $d^2$ has degree 172, but $-2(s+c_*)b_*$ has degree 180
and dominates it. Since $s$ and $c_*$ have the same leading component,
\[
E_0=-4(c_*)_0(b_*)_0,\qquad
(P_{96})_0=16(c_*)_0^2(b_*)_0^2.
\tag{42}
\]
The subtracted norm term has degree 358, below 360.
Thus the exact degree is \textbf{360}.
Its full leading component is
\[
2^{20}q_0^{184}w^{36}i^{20}Y^{76}
(\eta+\zeta)^{40}(q_0-F)^4,\qquad q_0=\beta J.
\tag{43}
\]
For example,
\[
[J^{184}w^{36}i^{20}Y^{76}\eta^{40}F^4]P_{96}
=2^{20}\beta^{184}>0.
\tag{44}
\]

\subsection{The direct-base equation of degree 332}

Replace $w$ by a positive input $x$ and set $X=q+x$.
Keep the same zero-offset main parameter and radicands. Use
\[
B=(qK_0+X)C-qF,\quad W=1+q(q-1)b,\quad
P_0=q^4(q-1),\quad M=P_0\mu^2H,
\]
\[
\begin{split}
Z={}&P_0\mu^2[H(n+V)+K_{\rm pos}]\\
&+H[\mu^2(q(q-1)Y-q^3WB-(q-1)X)-P_0e].
\end{split}
\tag{45}
\]
Evaluate the same unweighted norm (19)--(21).
The square gap and auxiliary index argument use $X\ge32$, not $q\mid X$,
so apply before the direct carrier recovers that divisibility.
The retained inequality $3\tau+f+X/q^4<HI<b/2$ supplies the corresponding
transport estimate.

For canonical completeness the direct parent has the exact map
\[
x=X-q,\qquad
n_{\rm direct}=n_{\rm scaled}+(q-1)^2bt_0+X/q^4>0.
\tag{46}
\]
Its numerator and denominator scale by $q$ and the norm output by $q^8$.
The late-divisibility proof and original source reconstruction remain valid
\cite{direct}.

The direct evaluator has 102 operations.
Its norm degrees are $(68,70,96,96,80)$, so the degree is 332.
The leading component is
\[
2^{20}q_0^{156}X_0^{36}i^{20}Y^{76}
(\eta+\zeta)^{40}(q_0-F)^4,\qquad X_0=q_0+x.
\tag{47}
\]
The coefficient of $J^{156}x^{36}i^{20}Y^{76}\eta^{40}F^4$ is
$2^{20}\beta^{156}>0$.
This companion trades six more operations for 28 fewer degrees.

\subsection{Retained intermediate equations}

The packet retains the weight-only reductions, before shortening the
transport carrier: 98 operations at degree 368 and 104 at degree 340.
They isolate the new radical-separation theorem.
The final 96/360 and 102/332 endpoints improve both degree and operations.

For the weight-only step the exact identity
\[
P_9-P_1
=16a(5a-s-b_*-c_*)(E_9+E_1)-512sc_*ab_*
\tag{48}
\]
has nonzero differences of exact degrees 362 and 334 respectively,
where $E_\lambda=d^2-2(s+c_*)(\lambda a+b_*)+(\lambda a-b_*)^2$.
Even this first step changes the polynomial, unlike the integer norm
rewrites.


\section{Merging sign certificates in the direct quartic}
\label{sec:guards}

The direct quartic uses a different source representation. Its arithmetic
body includes quotient/remainder rows, with separate nonnegativity
certificates for a remainder and its complementary slack. A shared positive
modulus lets two such certificates be replaced by one.

\begin{lemma}[Shared remainder bounds]
\label{lem:guard}
For integers $r,s,B$ with $r+s+1=B$ and $B\ge1$,
\[
(r\ge0\ \land\ s\ge0)\quad\longleftrightarrow\quad rs\ge0.
\tag{49}
\]
\end{lemma}

\begin{proof}
The forward implication is immediate. Conversely, nonnegative product
excludes opposite nonzero signs. Both values cannot be negative because
their sum is $B-1\ge0$. If either is zero, the other equals $B-1\ge0$.
\end{proof}

For integer variables, nonnegativity has the representation
\[
n=a^2+b^2+d(d+1).
\tag{50}
\]
The right side is nonnegative. Conversely, $4n+1$ for $n\ge0$ is not
of the excluded form $4^u(8v+7)$, so the three-square theorem represents
it by three squares. Exactly one coordinate is odd, by reduction
modulo four. Relabeling gives
$4n+1=(2a)^2+(2b)^2+(2d+1)^2$, proving (50).

Two linear nonnegative expressions normally require six private integer
witnesses. Their product instead uses three, with the quadratic row
$rs-a^2-b^2-d^2-d=0$.
The modulus sign is supplied by a retained anchor pair with the same
sum; no source interpretation is needed.

\subsection{The six literal pairs}

The indices below refer to the original arithmetic body, before projection.
Each anchor remains separately guarded, and its exact source rows prove
that the shared modulus is at least one.

\begin{table}[ht]
\centering
\begin{tabular}{@{}rrrr@{}}
\toprule
Remainder & Slack & Modulus & Retained anchor pair\\
\midrule
562&563&547&549,\ 550\\
585&586&574&576,\ 577\\
593&594&574&576,\ 577\\
601&602&574&576,\ 577\\
619&620&608&610,\ 611\\
637&638&626&628,\ 629\\
\bottomrule
\end{tabular}
\caption{Six applications of (49). The coefficient certificate checks
each $r+s+1-B=0$ identity by summing actual source rows.}
\label{tab:guards}
\end{table}

For example, raw rows 553--555 are
\[
-v_{585}-v_{586}+v_{589}=0,\quad
-1-v_{589}+v_{590}=0,\quad v_{574}-v_{590}=0.
\]
Their sum is $v_{574}-v_{585}-v_{586}-1=0$.
Three other unchanged rows give
$v_{574}-v_{576}-v_{577}-1=0$.
The retained guards on $v_{576}$ and $v_{577}$ establish $v_{574}\ge1$;
the new guard on $v_{585}v_{586}$ therefore restores both omitted signs.

The six replacements retain 63 old predicates and add six product
predicates, giving 69 instead of 75.
Every arithmetic coordinate and body equation is unchanged.
A new zero restores the twelve omitted nonnegative values by (49), and
fresh witnesses from (50) give the complete old system.
An old zero makes each product nonnegative, so fresh product witnesses
give the new system.
Thus their projections onto the arithmetic body coordinates are identical.

\subsection{One strict lower bound replaces two more signs}

There is a further exact implication in the same literal body:
\[
v_7=v_6^2,\ v_8=v_7^2,\ \ldots,\ v_{28}=v_{27}^2,\qquad
v_{28}=v_3+1.
\tag{51}
\]
There are 22 squaring rows, followed by the two original rows establishing
the last equality. On these integer equations,
\[
(v_3\ge0\ \land\ v_6\ge0)\quad\longleftrightarrow\quad v_6\ge1.
\tag{52}
\]
Indeed $v_6\ge1$ propagates through the squares to $v_{28}\ge1$,
giving $v_3\ge0$. Conversely $v_6=0$ would force $v_3=-1$,
contradicting the old guard; the old $v_6\ge0$ then implies $v_6\ge1$.
This covers the boundary $v_6=1,v_3=0$.

Replace the two old triples by the single guard
\[
v_6-1-a^2-b^2-d(d+1)=0.
\tag{53}
\]
The source now has 68 predicates.
A new zero first restores the 69-predicate system by (51)--(52);
the shared remainder proof then restores all 75 original predicates.
Only after those restorations is the C222 source interpretation applied.
Completeness uses fresh triples in the reverse order.
No omitted sign is justified by a semantic interpretation that itself
requires that sign.

\subsection{The complete 589-variable quartic}

The unprojected system has $1,259+3\cdot68=1,463$ integer coordinates and
$1,177+68=1,245$ quadratic rows.
The principal monic projection removes 874 coordinates and rows,
leaving 589 coordinates and 371 quadratic residuals.
Each deleted coordinate has a unique integral restoration, checked by
substituting in every remaining equation.
There is no rational relaxation or unrecorded free variable.

Let the delivered residuals be $f_0,\ldots,f_{370}$. The equation is
\[
\boxed{
Q_{589}(x_0,\ldots,x_{588})
=\sum_{j=0}^{370}f_j(x_0,\ldots,x_{588})^2=0,
\qquad x_i\in\Z.
}
\tag{54}
\]
Every coefficient of every residual is explicitly supplied in
\texttt{quartic/fused68/support8/equation.json}.
The fully collected polynomial is in the adjacent
\texttt{polynomial.json}, and the plain-text residual list gives the same
finite specification without a JSON reader.

Over $\Z$, (54) is equivalent to all residuals being zero.
Projection followed by guard restoration recovers the entire old
75-predicate source. Applying its ordinary C222/refutation calibration
therefore gives (1) with $D=\Z$ and $n=589$.

\begin{table}[ht]
\centering\small
\begin{tabular}{@{}lrrrr@{}}
\toprule
Profile & Variables & Rows & Monomials & Height\\
\midrule
Matched old source &604&372&9,676&2,445,790,974,704\\
Six pair replacements &592&372&8,167&2,445,790,974,704\\
Final principal &\textbf{589}&371&8,156&2,445,790,974,704\\
Final, quadratic first &595&377&7,595&6,956,012\\
Final, smaller coefficients &600&382&7,075&964,213\\
\bottomrule
\end{tabular}
\caption{Complete unrestricted-integer quartics, all of exact degree four.
The 600-variable alternative has every coefficient below one million.}
\label{tab:quartics}
\end{table}

The guard count alone would predict an eighteen-variable saving from the
six pairs. The actual projection saves twelve, because the products keep
six additional source coordinates alive. The strict bound then saves
three more variables. All reported counts are from complete equations
and checked inverses, rather than estimates.

\Needspace{13\baselineskip}
\subsection{The direct logical endpoint}

The unchanged arithmetic body directly encodes the numerical C222 calculus.
The named logical endpoint is the composition of
\[
\texttt{KEndpoint.c222\_iff\_zf\_refutation}
\]
and
\[
\texttt{ConstructibleChoice.zf\_refutation\_iff\_zfc\_refutation}.
\]
The first identifies C222 derivation of its falsehood code with finite
ZF refutation; the second transports refutability between ZF and ZFC.
Their exact immediate source files are included and pinned, but their
transitive Lean imports were not rebuilt here.
The ordinary arithmetic bridge and later guard reductions are explicitly
identified in \texttt{quartic/CALIBRATION\_DEPENDENCIES.md}.
The equation uses neither K215 simulation nor Jones' universal polynomial.

\section{Verification, reproducibility and proof boundaries}

\paragraph{Exact algebra.}
The norm checker multiplies all sixteen conjugate factors and compares
all 495 coefficients with the new schedules. For each of six delivered
integer variants, it proves equality of the complete old and new tails
over a free polynomial ring and authenticates every unchanged prefix.
The positive checker verifies the parity identities, octic norm,
complete register inventories, exact leading polynomials and degree claims.
The quartic checker verifies source-row identities, all monic integral
restorations, and every collected coefficient of the entire sum of squares.

\paragraph{Finite specification and provenance.}
The files contain every coefficient and every operation.
The BFS reconstructor decodes the entire 436-character seed, constructs
all 430 actions and 943 tiles, and reconstructs all four full integers.
Each raw payload matches the pinned complete-page value, rather than only
a modular residue or selected bits.
The integer packet includes all six authentic coefficient payloads,
and the direct quartic has literal coefficients with no parameter ports.

\paragraph{Unbounded mathematics.}
The square gap, radical rationality, index bounds, transport contradiction,
guard reconstruction and consistency calibration concern arbitrarily
large witnesses. The paper proofs supply those implications.
Modular controls are only supplementary transcription checks.
Neither a finite search finding no zeros nor large numerical experiments
would prove consistency.

\subsection{Claims that are not being made}

This packet supplies no new complete Lean or Rocq certificate, no encoded
object-language derivation of the entire calibration, and no proof of
$\Con$ itself. The source pins and retained licenses identify the inherited
dependencies. We did not rerun every historical native compiler or earlier
proof-assistant build in the repository.

The operation count is not a bit-complexity bound. The fixed native
coefficients are large, and legitimate witnesses may be much larger.
The direct quartic has smaller literal coefficients but many more variables.
No result is advertised as uniformly best in all those measures.

\subsection{Reproduction}

Run from the extracted packet:
\begin{lstlisting}
python3 verify_all.py
\end{lstlisting}
The exact checkers use the Python standard library. They check the supplied
finite artifacts without network access; optional builders are included too.

The report is rebuilt with:
\begin{lstlisting}
pdflatex -interaction=nonstopmode -halt-on-error report.tex
pdflatex -interaction=nonstopmode -halt-on-error report.tex
\end{lstlisting}
Recorded receipts and a checksum manifest accompany the files.

\section{What this round establishes}

The direct integer result saves fifteen actual variables and more than
fifteen percent of the expanded support at unchanged degree and coefficient
height. It uses evidence already present in the source: a common positive
remainder bound and a short squaring chain.

The compact positive result uses two genuine proof improvements.
The square-gap lemma replaces an asymmetric radical weight with a uniform
rationality argument. The auxiliary index bound then shows that the
existing witness already supplies enough growth for transport, eliminating
two additional arithmetic operations. The final scaled and direct equations
improve both degree and operation count over their respective parents.

The integer norm changes are algebraically simpler to certify.
Because the complete polynomials agree identically, there is no new
signed-domain soundness argument or exceptional specialization to resolve.
The packet retains both the symmetric derivation and the one-operation
cheaper schedule found in the pinned source.

Two tempting reductions were not promoted to theorems.
Replacing $\mu=q^2+1$ by $q+1$ destroys the first-index pinning inequality
used with $R<q^4$. Replacing the trace's additive positivity certificate
by a Pell guard can leave a zero-ordinate branch.
No delivered endpoint relies on either unproved replacement.

A substantially smaller variable count remains open. The present
constructions are concrete improvements with reproducible equations and
proofs, not evidence that a global minimum has been reached.

\appendix
\section{Complete fixed coefficients for the positive equations}
\label{app:constants}

This is the finite BFS compact44 recipe used by the pinned complete
equation page \cite{page}. It defines all four coefficients in (16)--(21)
without an external coefficient oracle. Every index and exponent is fixed
or bound in a finite sum. The script \texttt{reconstruct\_constants.py}
executes this recipe independently and compares all four entire raw byte
strings with the pinned values. The letter $K$ in this inherited recipe
denotes $K_0$ in the main text.

\begingroup
\small\setlength{\parskip}{4pt}
\subsection*{The four fixed integers}
Every index, set, sum, product and floor below defines a constant, not
another unknown. Put
\[
\rho=2^{67},\quad B=\rho^{8807},\quad\beta=B-1,\quad
G_m(x)=\sum_{j=0}^{m-1}x^j.
\]
Define
\[
\begin{split}
\mathcal E={}&\rg0{942}\cup\rg{983}{992}\cup\rg{1003}{1011}\\
&\cup\rg{1013}{1031}\cup\{1063,1064,1095,2195,3296\}.
\end{split}
\]
Then
\begin{align*}
M_C={}&B-1-\sum_{e\in\mathcal E\setminus\{1\}}\rho^e-2\rho^{1063},\\
M_F={}&B-1+4+(2^{51}-2)(\rho^{4393}+\rho^{5495})\\
&+\rho^{1013}G_{10}(\rho)+\rho^{1065}G_9(\rho)
  +\rho^{2195}+\rho^{3296}\\
&+(2^{65}-1)\rho^{8792}G_{13}(\rho)+(2^{19}-1)\rho^{8805},\\
K={}&\rho^{30}+\rho^{62}+\rho^{1100}+\rho^{2201}+\rho^{8849}\\
&+(\rho^{4393}+\rho^{5495})\Phi(1/\rho),
\end{align*}
where
\[
\begin{split}
\Phi(x)={}&(1+2^{51})G_{943}(x)
 -2^{10}\sum_{t\in\mathcal T}v(t)x^{\varrho(t)}\\
&+\sum_{(b,e,m)\in\mathcal C}2^b x^e G_m(2x),\\
\mathcal C={}&\{(10,1023,8),(18,1013,9),(27,1003,8),\\
&(35,983,9),(44,1031,1),(45,1022,1),(46,1011,1),\\
&(47,992,1),(48,1095,1),(49,2195,1),(50,3296,1)\}.
\end{split}
\]
The rational intermediate in K cancels exactly; all four results are positive integers.

\subsection*{The 943 fixed tiles and their ranks}
Use the actions $u_j,v_j,d_j$ decoded below. Let
\[
\begin{gathered}
h(q,b)=2(q+1)+b,\quad
\mathcal R=\{u_j:d_j=1,\ 0\le j<430\},\\
\mathcal L=\{u_j:d_j=0,\ 0\le j<430\},\\
r_q=6+|\mathcal R\cap\rg0{q-1}|,\quad
l_q=205+|\mathcal L\cap\rg0{q-1}|.
\end{gathered}
\]
An interval ending at $-1$ is empty. The following union is $\mathcal T$;
the index ranges in each row are fixed:
\[
\setlength{\arraycolsep}{1.4pt}\begin{array}{c|l}
\text{range}&\text{tile }(w,n,e,s)\\\hline
j=2q+b<430&((1-d_j)l_{u_j},h(q,b),d_jr_{u_j},v_j)\\
b=0,1&(0,h(215,b),0,h(215,b)),\ (0,b,0,b)\\
q\in\mathcal R,\ b=0,1&(r_q,b,0,h(q,b))\\
q\in\mathcal L,\ b=0,1&(0,b,l_q,h(q,b))\\
t\in\{0,1,432,433\}&(0,t,0,438)\\
&(2,434,2,0),\ (2,435,3,1)\\
&(3,436,4,h(u_0,0)),\ (4,436,4,0),\ (5,0,0,0)\\
t=434,435,436&(1,438,1,t)\\
a\in\{0,1,4\},\ b=0,1,2&(a,437,b,437)
\end{array}
\]
Put $s_*=(2,435,3,1)$, $f_*=(5,0,0,0)$, and $\mathcal T'=\mathcal T\setminus\{s_*,f_*\}$.
For $t=(w,n,e,s)$ set
\[
\begin{gathered}
\kappa(t)=2^{30}w+2^{20}n+2^{10}e+s,\\
v(t)=w+2^8n+2^{17}e+2^{25}s+127\,2^{34},\\
\varrho(s_*)=0,\quad\varrho(f_*)=1,\\
\varrho(t)=2+\sum_{z\in\mathcal T'}
 \fl{\frac{2^{40}+\kappa(t)-\kappa(z)-1}{2^{40}}}
 \quad(t\in\mathcal T').
\end{gathered}
\]

\subsection*{A closed decoder for the 430 actions}
Put $D_b(A,k)=\fl{A/b^k}\bmod b$, where
$a\bmod b=a-b\fl{a/b}$. For $0\le q<215$, $b=0,1$, $j=2q+b$:
\[
\begin{gathered}
N=\mathcal A\bmod2^{636},\quad
T_0=\fl{\mathcal A/2^{636}}\bmod2^{1675},\\
E_0=\fl{\mathcal A/2^{2311}},\\
\nu_j=D_2(N,j),\quad \sigma_j=\sum_{k=0}^{j-1}\nu_k,\\
m_q=D_2(N,430+q),\quad \tau_q=\sum_{k=0}^{q-1}m_k,\\
F_q=14(1-m_q)+m_qD_{16}(E_0,\tau_q),\\
u_j=\nu_j(1+\sigma_j)+(1-\nu_j)D_{216}(T_0,j-\sigma_j),\\
v_j=D_2(F_q,2b),\quad d_j=D_2(F_q,2b+1).
\end{gathered}
\]
Empty sums are zero. Here $(v_0,d_0,u_0)=(1,1,1)$; the constructor tile
uses this actual first target, including its state renaming.

\subsection*{The fixed seed integer}
The alphabet, in order of digit values $0,\ldots,61$, is
{\fontsize{7.7}{8.2}\selectfont\ttfamily 0123456789ABCDEFGHIJKLMNOPQRSTUVWXYZabcdefghijklmnopqrstuvwxyz}.\par
Concatenate the following 436 characters as digits $e_0,\ldots,e_{435}$;
put $\mathcal A=\sum_{j=0}^{435}e_j62^{435-j}$.\par
\noindent\begin{minipage}{\linewidth}
\fontsize{7.7}{8.2}\selectfont\ttfamily\setlength{\parskip}{0pt}
BxEnhvy6NipzkWmROo0GsA3uK3twWLxscLq20UTumy5Cdu461LwD6g7RDNzF8Y\par
iYDFTpB8ehAypd7d6OsOL4APamnJeNXQypUezCmW9n4oXP6XfgmiWuAWZ8ZJ7U\par
OP7gu4kVdlGZqZHULQMar1dXNeJW4i7oNVI8usp8BlovjjZbwPCovZI328fRUq\par
4q92R81s5GjWQaQfGhkq9UYMK3rRMaBmsOIBkJqc2cj2VRIZ6mjzkBIOkW9SjZ\par
FzeA6J5u4j3sBL3SdgwXzM9mHvQivu8EJ1TBTpRhHDqZIbmYgwQ8VqTAFotl6F\par
We1YC554GKrVHsA8oUGKSA6F6k5KREBSD0Tgd7bFALY84LLbpw574BQyvdoHrl\par
0Oe35uwDf15DpT6VJWZBjSbHPpPAINYTte0eUxnTFTKBlGj8Am6sb1gLldh8WZ\par
8P\par
\end{minipage}

\endgroup

\clearpage
\section{Delivered file organization}

\begin{longtable}{@{}p{0.31\textwidth}p{0.63\textwidth}@{}}
\toprule
Path & Contents\\
\midrule
\endhead
\texttt{report.tex}, \texttt{report.pdf}&
The report, with its companion LaTeX input files.\\
\texttt{positive/}&
Final 96/360 and 102/332 circuits; weight-only checkpoints; mathematical
proofs, exact checks, original coefficients and pinned source proofs.\\
\texttt{positive/coefficients\_bfs/}&
The four completely reconstructed coefficients used by the main displayed
equation, with a binding for the BFS circuit variants.\\
\texttt{integer/}&
Preferred 126-, 138-, and 137-operation integer circuits, symmetric
alternatives, original parents, all six coefficients and exact norm checks.\\
\texttt{quartic/}&
All quadratic residuals, complete quartic coefficients, source predicates,
projection certificates, builders, exact checkers and matched parent profiles.\\
\texttt{reconstruct\_constants.py}&
Finite reconstruction from the printed seed, actions and tiles.\\
\texttt{verify\_all.py}&
Runs the selected exact checks and records their outcomes.\\
\texttt{PROVENANCE.json}&
Frozen repository snapshot, source references and claim boundaries.\\
\texttt{SHA256SUMS.txt}&
Hashes of package files, excluding this checksum file itself.\\
\bottomrule
\end{longtable}

\begin{thebibliography}{99}
\small\raggedright
\bibitem{zero}
ProveIt contributors. \emph{Ten-positive zero-offset equation}.
\url{https://github.com/VladimirReshetnikov/ProveIt/tree/84e16e43092614bae262ca59651cb9d77e31a255/Computability/DiophantineConsistency/candidates/ten-positive-zero-offset}.
\bibitem{direct}
ProveIt contributors. \emph{Direct positive base with the zero-offset main parameter}.
\url{https://github.com/VladimirReshetnikov/ProveIt/tree/84e16e43092614bae262ca59651cb9d77e31a255/Computability/DiophantineConsistency/candidates/ten-positive-direct-zero-offset}.
\bibitem{trace}
ProveIt contributors. \emph{Thirteen-integer quadratic trace equation}.
\url{https://github.com/VladimirReshetnikov/ProveIt/tree/84e16e43092614bae262ca59651cb9d77e31a255/Computability/DiophantineConsistency/candidates/integer13-quadratic-trace}.
\bibitem{scalar}
ProveIt contributors. \emph{Scalar carrier in twelve unrestricted integers}.
\url{https://github.com/VladimirReshetnikov/ProveIt/tree/84e16e43092614bae262ca59651cb9d77e31a255/Computability/DiophantineConsistency/candidates/scalar-carrier12}.
\bibitem{quartic75}
ProveIt contributors. \emph{Complete 75-predicate direct quartic family}.
\url{https://github.com/VladimirReshetnikov/ProveIt/tree/84e16e43092614bae262ca59651cb9d77e31a255/Computability/DiophantineConsistency/research/factorial-radius-output-sign-profile}.
\bibitem{page}
ProveIt contributors. \emph{Complete ten-positive equation page and coefficients}.
\url{https://github.com/VladimirReshetnikov/ProveIt/tree/84e16e43092614bae262ca59651cb9d77e31a255/Computability/DiophantineConsistency/candidates/one-page-short10}.
\bibitem{sun}
Zhi-Wei Sun. \emph{Ten unknowns for Hilbert's tenth problem over the integers}.
arXiv:2609.01594, 2026. \url{https://arxiv.org/abs/2609.01594}.
\end{thebibliography}


\end{document}

